\section{Korean}

The Korean tradition preserves an old layer of creation lore in the songs of its village and island shamans, told outside the later Buddhist and Confucian frames. The fullest account is the \textbf{Cheonjiwang Bonpuri}, the origin-song of the heavenly king, sung by the shamans of Jeju Island. Alongside it runs an older mainland strand of a giant grandmother creator, \textbf{Mago}, who shaped the land itself.

\subsection{The Creation Model}

At first heaven and earth were one fused mass with no gap between them. Then the mass split. The light rose to become the sky and the heavy settled to become the ground, and the gap between them opened the room of the world. But the early world was unlivable. Two suns blazed in the day and burned the living, and two moons froze the night.

The king of heaven came down and fathered twin sons by an earthly woman. The two brothers climbed to the sky and took up iron bows. One shot down the extra sun, the other shot down the extra moon, leaving a single sun and a single moon and a world fit to live in. Then the brothers held a contest to decide who would rule the living and who the dead. The younger won the world of the living by a trick, which is why the human world still holds disorder and wrong, while the elder took the world of the dead and rules it with strict justice.

\subsection{Key Motifs}

The heavenly king is \textbf{Cheonjiwang}. His twin sons are the Great Star King (\textbf{Daebyeol-wang}) and the Little Star King (\textbf{Sobyeol-wang}). The undivided beginning and its opening is the loss and gain of the \textbf{gawp}, the divide between the realms. The flower-growing contest decides the split rule of the living world and the world of the dead. In the separate mainland tradition the grandmother giant \textbf{Mago}, also called \textbf{Magohalmi}, forms the mountains and rivers while the male powers set out heaven, earth, sun, and moon.

This tradition states the shared pattern as the opening of one fused mass into the paired realms of sky and earth, and then the correction of an overflowing world down to a single ordered sun and moon. The going-out is the splitting and the doubling, and the return toward balance is the shooting-down that leaves one of each.
